\documentclass[UTF8,12pt,fontset=windows]{ctexart}
\usepackage[paperwidth=240mm,paperheight=200mm,left=14mm,right=14mm,top=9mm,bottom=12mm]{geometry}
\usepackage{amsmath,amssymb,mathtools,bm,graphicx,xcolor}
\pagestyle{empty}
\setlength{\parindent}{0pt}
\setlength{\parskip}{4pt}
\setlength{\emergencystretch}{2em}
\xeCJKDeclareCharClass{CJK}{"2460 -> "24FF}
\xeCJKDeclareCharClass{CJK}{"2160 -> "2188}
\begin{document}
{\large\bfseries 2002.1}\quad {\small 原卷一（1）}

设函数 $f(x)=\begin{cases}\dfrac{1-\mathrm{e}^{\tan x}}{\arcsin\dfrac{x}{2}}, & x>0,\\[4pt] a\mathrm{e}^{2x}, & x\le0,\end{cases}$ 在 $x=0$ 处连续，则 $a=$\underline{\hspace{20mm}}。

\vfill
\newpage
{\large\bfseries 2002.2}\quad {\small 原卷一（2）}

位于曲线 $y=x\mathrm{e}^{-x}\ (0\le x<+\infty)$ 下方，$x$ 轴上方的无界图形的面积是\underline{\hspace{20mm}}。

\vfill
\newpage
{\large\bfseries 2002.3}\quad {\small 原卷一（3）}

微分方程 $yy''+y'^2=0$ 满足初始条件 $y\big|_{x=0}=1,\ y'\big|_{x=0}=\dfrac{1}{2}$ 的特解是\underline{\hspace{20mm}}。

\vfill
\newpage
{\large\bfseries 2002.4}\quad {\small 原卷一（4）}

$\lim_{n\to\infty}\dfrac{1}{n}\left(\sqrt{1+\cos\dfrac{\pi}{n}}+\sqrt{1+\cos\dfrac{2\pi}{n}}+\cdots+\sqrt{1+\cos\dfrac{n\pi}{n}}\right)=$\underline{\hspace{20mm}}。

\vfill
\newpage
{\large\bfseries 2002.5}\quad {\small 原卷一（5）}

矩阵 $\begin{pmatrix}0&-2&-2\\2&2&-2\\-2&-2&2\end{pmatrix}$ 的非零特征值是\underline{\hspace{20mm}}。

\vfill
\newpage
{\large\bfseries 2002.6}\quad {\small 原卷二（1）}

设函数 $f(u)$ 可导，$y=f(x^2)$。当自变量 $x$ 在 $x=-1$ 处取得增量 $\Delta x=-0.1$ 时，相应的函数增量 $\Delta y$ 的线性主部为 $0.1$，则 $f'(1)=(\quad)$。

(A) $-1$.  (B) $0.1$.  (C) $1$.  (D) $0.5$.

\vfill
\newpage
{\large\bfseries 2002.7}\quad {\small 原卷二（2）}

设函数 $f(x)$ 连续，则下列函数中，必为偶函数的是 $(\quad)$。

(A) $\int_0^x f(t^2)\,\mathrm{d}t$.  (B) $\int_0^x f^2(t)\,\mathrm{d}t$.  (C) $\int_0^x t[f(t)-f(-t)]\,\mathrm{d}t$.  (D) $\int_0^x t[f(t)+f(-t)]\,\mathrm{d}t$.

\vfill
\newpage
{\large\bfseries 2002.8}\quad {\small 原卷二（3）}

设 $y=y(x)$ 是二阶常系数微分方程 $y''+py'+qy=\mathrm{e}^{3x}$ 满足初始条件 $y(0)=y'(0)=0$ 的特解，则当 $x\to0$ 时，函数 $\dfrac{\ln(1+x^2)}{y(x)}$ 的极限 $(\quad)$。

(A) 不存在.  (B) 等于1.  (C) 等于2.  (D) 等于3.

\vfill
\newpage
{\large\bfseries 2002.9}\quad {\small 原卷二（4）}

设函数 $y=f(x)$ 在 $(0,+\infty)$ 内有界且可导，则 $(\quad)$。

(A) 当 $\lim_{x\to+\infty}f(x)=0$ 时，必有 $\lim_{x\to+\infty}f'(x)=0$.

(B) 当 $\lim_{x\to+\infty}f'(x)$ 存在时，必有 $\lim_{x\to+\infty}f'(x)=0$.

(C) 当 $\lim_{x\to0^+}f(x)=0$ 时，必有 $\lim_{x\to0^+}f'(x)=0$.

(D) 当 $\lim_{x\to0^+}f'(x)$ 存在时，必有 $\lim_{x\to0^+}f'(x)=0$.

\vfill
\newpage
{\large\bfseries 2002.10}\quad {\small 原卷二（5）}

设向量组 $\alpha_1,\alpha_2,\alpha_3$ 线性无关，向量 $\beta_1$ 可由 $\alpha_1,\alpha_2,\alpha_3$ 线性表示，而向量 $\beta_2$ 不能由 $\alpha_1,\alpha_2,\alpha_3$ 线性表示，则对于任意常数 $k$，必有 $(\quad)$。

(A) $\alpha_1,\alpha_2,\alpha_3,k\beta_1+\beta_2$ 线性无关.

(B) $\alpha_1,\alpha_2,\alpha_3,k\beta_1+\beta_2$ 线性相关.

(C) $\alpha_1,\alpha_2,\alpha_3,\beta_1+k\beta_2$ 线性无关.

(D) $\alpha_1,\alpha_2,\alpha_3,\beta_1+k\beta_2$ 线性相关.

\vfill
\newpage
{\large\bfseries 2002.11}\quad {\small 原卷三}

已知曲线的极坐标方程是 $r=1-\cos\theta$，求该曲线上对应于 $\theta=\dfrac{\pi}{6}$ 处的切线与法线的直角坐标方程。

\vfill
\newpage
{\large\bfseries 2002.12}\quad {\small 原卷四}

设 $f(x)=\begin{cases}2x+\dfrac{3}{2}x^2, & -1\le x<0,\\[4pt] \dfrac{x\mathrm{e}^x}{(\mathrm{e}^x+1)^2}, & 0\le x\le1,\end{cases}$ 求函数 $F(x)=\int_{-1}^x f(t)\,\mathrm{d}t$ 的表达式。

\vfill
\newpage
{\large\bfseries 2002.13}\quad {\small 原卷五}

已知函数 $f(x)$ 在 $(0,+\infty)$ 内可导，$f(x)>0$，$\lim_{x\to+\infty}f(x)=1$，且满足

$$\lim_{h\to0}\left[\frac{f(x+hx)}{f(x)}\right]^{\frac{1}{h}}=\mathrm{e}^{\frac{1}{x}},$$

求 $f(x)$。

\vfill
\newpage
{\large\bfseries 2002.14}\quad {\small 原卷六}

求微分方程 $x\,\mathrm{d}y+(x-2y)\,\mathrm{d}x=0$ 的一个解 $y=y(x)$，使得由曲线 $y=y(x)$ 与直线 $x=1,\ x=2$ 以及 $x$ 轴所围成的平面图形绕 $x$ 轴旋转一周的旋转体体积最小。

\vfill
\newpage
{\large\bfseries 2002.15}\quad {\small 原卷七}

某闸门的形状与大小如图所示，其中直线 $l$ 为对称轴，闸门的上部为矩形 $ABCD$，下部由二次抛物线与线段 $AB$ 所围成。当水面与闸门的上端相平时，欲使闸门矩形部分承受的水压力与闸门下部承受的水压力之比为 $5:4$，闸门矩形部分的高 $h$ 应为多少 $\mathrm{m}$（米）？

\begin{center}\includegraphics[width=.62\textwidth,height=65mm,keepaspectratio]{_题图/2002.15.png}\end{center}

\vfill
\newpage
{\large\bfseries 2002.16}\quad {\small 原卷八}

设 $0<x_1<3$，$x_{n+1}=\sqrt{x_n(3-x_n)}\ (n=1,2,\cdots)$，证明数列 $\{x_n\}$ 的极限存在，并求此极限。

\vfill
\newpage
{\large\bfseries 2002.17}\quad {\small 原卷九}

设 $0<a<b$，证明不等式

$$\frac{2a}{a^2+b^2}<\frac{\ln b-\ln a}{b-a}<\frac{1}{\sqrt{ab}}.$$

\vfill
\newpage
{\large\bfseries 2002.18}\quad {\small 原卷十}

设函数 $f(x)$ 在 $x=0$ 的某邻域内具有二阶连续导数，且 $f(0)\neq0$，$f'(0)\neq0$，$f''(0)\neq0$。证明：存在唯一的一组实数 $\lambda_1,\lambda_2,\lambda_3$，使得当 $h\to0$ 时，$\lambda_1f(h)+\lambda_2f(2h)+\lambda_3f(3h)-f(0)$ 是比 $h^2$ 高阶的无穷小。

\vfill
\newpage
{\large\bfseries 2002.19}\quad {\small 原卷十一}

已知 $A,B$ 为3阶矩阵，且满足 $2A^{-1}B=B-4E$，其中 $E$ 是3阶单位矩阵。

(1) 证明：矩阵 $A-2E$ 可逆；

(2) 若 $B=\begin{pmatrix}1&-2&0\\1&2&0\\0&0&2\end{pmatrix}$，求矩阵 $A$。

\vfill
\newpage
{\large\bfseries 2002.20}\quad {\small 原卷十二}

已知4阶方阵 $A=(\alpha_1,\alpha_2,\alpha_3,\alpha_4)$，$\alpha_1,\alpha_2,\alpha_3,\alpha_4$ 均为4维列向量，其中 $\alpha_2,\alpha_3,\alpha_4$ 线性无关，$\alpha_1=2\alpha_2-\alpha_3$。如果 $\beta=\alpha_1+\alpha_2+\alpha_3+\alpha_4$，求线性方程组 $Ax=\beta$ 的通解。

\vfill
\end{document}
